\begin{afterword}
\summaryhead

The post argues that religion is a symptom of a wider failure of sanity. If every trace of religion vanished tomorrow, the failures that made it possible would remain. Religion is like a dead canary in a coal mine: it shows that something is wrong with the air.

So rationalists should also teach general methods of reasoning, never mentioning religion, until religion ``goes underwater''. The author proposes a test: teach such a course, never hint at religion, and five years later compare how many students have lost their religion with a control group. A list of the author's earlier topics is offered as material.

The second argument is that a religious scientist is not applying any explicit understanding of evidence, Occam's Razor and the like, although an introduction to all of it would fit in one undergraduate course. Nobel laureates such as Richard Smalley and Robert Aumann are religious, and their colleagues did not correct them, so the ``sanity waterline'' is very low even in science. Even a complete victory of the neo-atheists would leave the real work of rationalists only beginning.

\responsehead

The post's proposal is its strongest part. It turns a slogan into a design that could be tested: a course that never mentions religion, a control group, and a count five years later. The next day's post, ``A Sense That More Is Possible'', asks for the same kind of long, controlled follow-up of rationality training.

The premise the test would check, that general methods taught without mentioning religion would lower religious belief, is offered as a hope; no evidence is given for it yet. The material proposed is a list of the author's own posts, and the post itself grants that one of them, Dark Side Epistemology, would be hard to teach without religion.

The second argument claims more than its evidence shows. From the bare fact that a scientist is religious, the post concludes, as something ``We now know'', that the person is applying none of a list of skills. This is the diagnosis of ``Outside the Laboratory'', which the post links, now without that post's ``to me'' and ``probably''. As there, no real scientist's reasoning is examined.

For Aumann, one of the two laureates the post names, whose own account of his religion differs from the post's picture, the notes on ``Crisis of Faith'' give the details.

In short: a testable proposal whose premise, that general methods would dissolve religion, is what the test would check, together with a diagnosis of religious scientists drawn from the bare fact of their belief.
\end{afterword}
